% !TEX root = 03.tex
\documentclass[a4paper,10pt]{article}

\usepackage{pdfsync}
\usepackage[utf8]{inputenc}
\usepackage{microtype}
\usepackage{mathtools}
\usepackage{amsthm,amsfonts,amsmath,amssymb}
\usepackage{mathabx}
% \usepackage{MnSymbol}
\usepackage{hyperref}
% \usepackage{bbold}
\usepackage[textwidth=2cm,textsize=small]{todonotes}
\usepackage{subcaption}
\usepackage{enumitem}
\usepackage{xspace}
% \usepackage{stmaryrd}
\usepackage{verbatim}
\usepackage{cleveref}
\usepackage{csquotes}
\usepackage{cite}

\bibliographystyle{plain}

\hypersetup{
    colorlinks,
    linkcolor={red!50!black},
    citecolor={blue!50!black},
    urlcolor={blue!80!black}
}

\author{}

\newcommand{\ignore}[1]{}

\newcommand{\eg}{\textit{e.g.}\xspace}
\newcommand{\ie}{\textit{i.e.}\xspace}
\newcommand{\ea}{\textit{et al.}\xspace}
\newcommand{\etc}{\textit{etc.}\xspace}
\newcommand{\resp}{resp.\xspace}
\newcommand{\st}{s.t.\xspace}

\newcommand{\sep}{\;|\;}
\renewcommand{\rule}[2]{\dfrac{#1}{#2}}
\newcommand{\set}[1]{\left\{#1\right\}}
\newcommand{\setof}[2]{\set{#1 \;\middle|\; #2}}
\newcommand{\tuple}[1]{\left( #1 \right)}
\newcommand{\pair}[2]{\tuple {#1, #2}}
\newcommand{\closure}[3]{\tuple {#1, #2, #3}}
\newcommand{\sem}[1]{\llbracket{#1}\rrbracket}
\newcommand{\powerset}[1]{\mathcal P(#1)}
\newcommand{\map}[2]{\mathsf{map}\; #1\; #2}
\newcommand{\filter}[2]{\mathsf{filter}\; #1\; #2}
\newcommand{\floor}[1]{\left\lfloor #1 \right\rfloor}
\newcommand{\upwardclosure}[1]{#1\uparrow}
\newcommand{\card}[1]{\left| #1 \right|}
\newcommand{\lang}[1]{L(#1)}
\newcommand{\goesto}[1]{\xrightarrow{#1}}
\newcommand{\e}{\varepsilon}
\newcommand{\supp}[1]{\mathsf{supp}(#1)}
\newcommand{\poly}[2]{#1[#2]}
\newcommand{\rat}[2]{#1(#2)}
\newcommand{\powerseries}[2]{#1[\![#2]\!]}
\newcommand{\dollar}{\$}
\newcommand{\coef}[2]{[#1]#2}

\newcommand{\C}{\mathbb C}
\newcommand{\N}{\mathbb N}
\newcommand{\Z}{\mathbb Z}
\newcommand{\Q}{\mathbb Q}
\newcommand{\K}{\mathbb K}
\newcommand{\R}{\mathbb R}

\newcommand{\NL}{\textsf{NL}}
\newcommand{\NC}[1]{$\textsf{NC}^{#1}$}
\newcommand{\PTIME}{\textsf{PTIME}}
\newcommand{\PSPACE}{\textsf{PSPACE}}

\newtheorem{exercise}{Exercise}

\newenvironment{solution}{\color{teal} \begin{proof}[\protect{\textnormal{\emph{Solution.}}}]}{\end{proof}}

\let\oldqedhere\qedhere

\newif\ifstartedinmathmode
\renewcommand*{\qedhere}{%
  \relax\ifmmode\startedinmathmodetrue\else\startedinmathmodefalse\fi%
  \ifstartedinmathmode\tag*{\protect\oldqedhere}\else\ifinlist{\hfill\oldqedhere}{\oldqedhere}\fi%
}

\crefname{equation}{}{}
\crefname{exercise}{Exercise}{Exercises}


\title{Formal power series and applications - Tutorial 03}
\date{2026/10/20}
\author{Yahia Idriss Benalioua \\ \small University of Warsaw}

\begin{document}

\maketitle

\section*{Hadamard product}

Let $x = \tuple{x_1, \dots, x_d}$ be a tuple of commuting indeterminates.
%
The \emph{Hadamard product} of two multivariate power series $f, g \in \powerseries \Q x$ is
%
\begin{align*}
    f \hadamard g = \sum_{n \in \N^d} f_n g_n \cdot x^n.
\end{align*}

\begin{exercise}
    Show that the Hadamard product of two univariate rational power series is rational.
\end{exercise}

\begin{exercise}
    Is the Hadamard product of two multivariate rational power series necessarily rational?
\end{exercise}

\section*{D-finite power series}

\begin{exercise}
    Consider power series $f_1, \dots, f_k \in \powerseries \Q x$ satisfying a \emph{system of linear differential equations} over $\rat \Q x$
    %
    \begin{align}
        f_i' = A_{i1} \cdot f_1 + \dots + A_{ik} \cdot f_k,
        \quad A_{ij} \in \rat \Q x,
        \quad 1 \leq i \leq k.
    \end{align}
    %
    Show that $f_1$ is D-finite.
\end{exercise}

\begin{exercise}
    Are D-finite power series closed under composition?
\end{exercise}

\begin{exercise}
    Let $f \in \powerseries \Q {x_1, \dots, x_k}$ be a multivariate D-finite power series
    and let $g_1, \dots, g_k \in \powerseries \Q {y_1, \dots, y_\ell}$ be algebraic multivariate power series without constant term,
    $g_1(0) = \dots = g_k(0) = 0$.
    %
    Show that $f(g_1, \dots, g_k)$ is D-finite.
\end{exercise}

\section*{Unambiguous Parikh automata}

\begin{exercise}
  The \emph{ambiguity problem} for Parikh automata takes as input a Parikh automaton
  and asks whether it is unambiguous.
  Is the ambiguity problem decidable?
\end{exercise}

\bibliography{../_bibliography/references}

\end{document}
